\documentclass[webpdf,contemporary,large,namedate]{oup-authoring-template}%

%\PassOptionsToPackage{hyphens}{url}
%\PassOptionsToPackage{colorlinks,linkcolor=blue,urlcolor=blue,citecolor=blue,anchorcolor=blue}{hyperref}

\DeclareMathOperator*{\argmax}{argmax}

\usepackage{lmodern}
\usepackage{algorithm}
\usepackage[noEnd=true,commentColor=darkgray]{algpseudocodex}
\tikzset{algpxIndentLine/.style={draw=darkgray}}
\usepackage{setspace}
\renewcommand{\ttdefault}{cmtt}

\begin{document}
\journaltitle{TBD}
\DOI{TBD}
\copyrightyear{2026}
\pubyear{2026}
\access{Advance Access Publication Date: Day Month Year}
\appnotes{Preprint}
\firstpage{1}

\title[Fast read alignment with minibwa]{Fast genomic read alignment with minibwa}
\author[1,2,3,$\ast$]{Heng Li\ORCID{0000-0003-4874-2874}}
\author[4]{Nils Homer\ORCID{0009-0007-7860-1155}}
\address[1]{Department of Data Science, Dana-Farber Cancer Institute, 450 Brookline Ave, Boston, MA 02215, USA}
\address[2]{Department of Biomedical Informatics, Harvard Medical School, 10 Shattuck St, Boston, MA 02215, USA}
\address[3]{Broad Insitute of MIT and Harvard, 415 Main St, Cambridge, MA 02142, USA}
\address[4]{Fulcrum Genomics LLC, 240 Elm Street, Somerville, MA, 02144, USA}
\corresp[$\ast$]{Corresponding author. \href{mailto:hli@ds.dfci.harvard.edu}{hli@ds.dfci.harvard.edu}}

%\received{Date}{0}{Year}
%\revised{Date}{0}{Year}
%\accepted{Date}{0}{Year}

\abstract{
\sffamily\footnotesize
\textbf{Motivation:}
BWA-MEM remains a popular short-read mapper especially for the purpose of variant calling.
Several groups have accelerated this algorithm as it has been the performance bottleneck of many current workflows.
However, constrained by the original design,
these drop-in replacements could only achieve limited speedup.
Breaking changes to BWA-MEM are required for further improvement.\vspace{0.5em}\\
\textbf{Results:}
We developed minibwa for aligning short and accurate long reads against a reference genome.
It combines BWA-MEM variable-length seeding with minimap2 chaining and base alignment.
It speeds up BWA-MEM2 further with additional prefetch for seeding,
new heuristics to skip unnecessary mate rescue and reduced effort in highly repetitive regions
where reads would anyway be wrongly mapped due to structural changes.
Minibwa is over three times as fast as BWA-MEM and over twice as fast as BWA-MEM2 at comparable accuracy.
It also natively supports directional bisulfite sequencing data.
\vspace{0.5em}\\
\textbf{Availability and implementation:}
\url{https://github.com/lh3/minibwa}
}

\maketitle

\section{Introduction}

Read alignment is indispensable to most reference-based sequence analyses such as variant calling and methylation profiling.
Although over 100 read mappers have been published~\citep{Alser:2021aa,DBLP:journals/csur/ProusalisGKNKAPSG26},
only several are maintained and remain widely used.
The two most notable short-read mappers for DNA sequencing data are Bowtie2~\citep{Langmead:2012fk} and BWA-MEM~\citep{Li:2013aa}, both based on FM-index~\citep{DBLP:conf/focs/FerraginaM00}.
While Bowtie2 is as popular for general purposes, BWA-MEM is more often used for variant calling.

BWA-MEM is the sole performance and cost bottleneck in variant calling pipelines.
Given that human Whole-Genome Sequencing (WGS) accounts for the vast majority of sequencing data volume, this inefficiency is particularly pronounced.
There have been several efforts to accelerate the original BWA-MEM algorithm.
BWA-MEM2~\citep{DBLP:conf/ipps/VasimuddinMLA19} revamps the memory layout of FM-index
and accelerates extension alignment with Single Instruction Multiple Data (SIMD) based algorithm.
It is typically 50--100\% faster and outputs identical alignment.
BWA-MEME~\citep{Jung:2022aa} further speeds up BWA-MEM2 with learned index and still promises identical output.
Sentieon BWA is a proprietary drop-in replacement of BWA-MEM.
Parabricks~\citep{OConnell:2023aa,Zhu2025.07.23.666378} provides a GPU-based partial reimplementation of BWA-MEM
that produces alignment functionally equivalent~\citep{Regier:2018aa} to BWA-MEM alignment.

However, designed more than a decade ago, BWA-MEM has several limitations that prevent further optimization but can be alleviated if we accept breaking changes.
First, BWA-MEM uses a complex algorithm to seed alignment~\citep{Li:2012fk}.
It is challenging to apply the latency hiding technique, which we will explain later, to speed up this step.
Second, BWA-MEM chaining is heuristic and inferior to the minimap2 chaining algorithm~\citep{Li:2018ab}.
It likely does not work well for short reads longer than 500bp, which will become more common with Roche's SBX sequencing technology~\citep{Kokoris2025.02.19.639056}.
The boundary between short and long reads is blurring.
Third, BWA-MEM uses a non-SIMD-based extension algorithm with intricate boundary conditions.
This greatly complicates SIMD-based equivalence in BWA-MEM2 and increases maintenance burden.
Fourth, BWA-MEM tries too hard in highly repetitive regions.
Due to rapid evolution in centromeres~\citep{Gao2025.12.09.693231}, 
exact base alignment in centromeres is challenging even with complete assemblies~\citep{Bzikadze:2022aa}.
Exhaustive alignment in centromeres is a waste of effort
as we cannot accurately place centromeric short reads anyway due to large structural differences.

In addition to algorithms for aligning standard DNA sequencing data,
there are also many algorithms for aligning bisulfite sequencing (BS-seq) reads.
These include
BSMAP~\citep{Xi:2009aa},
Bismark~\citep{Krueger:2011aa},
BS-Seeker~\citep{Chen:2010aa,Guo:2013aa,Huang:2018aa},
BWA-Meth~\citep{Pedersen:2014aa},
BitMapperBS~\citep{Cheng:2018aa},
gemBS~\citep{Merkel:2019aa},
HISAT-3N~\citep{Zhang:2021ac},
BSBolt~\citep{Farrell:2021aa},
abismal~\citep{de_Sena_Brandine_2021},
FAME~\citep{Fischer_2023},
BISCUIT~\citep{Zhou:2024ac},
RMapAlign3N~\citep{Muller:2025aa}
and ARYANA-BS~\citep{Nikaein_2025}.
Several generic short-read aligners, such as
GSNAP~\citep{Wu:2010uq},
segemehl~\citep{Otto:2014aa}
and Arioc~\citep{Wilton:2018aa},
also natively support BS-seq.

Among these algorithms, Bismark~\citep{Krueger:2011aa} is by far the most cited,
partly because it is a wrapper around stock Bowtie2~\citep{Langmead:2012fk}
which is battle tested and whose behavior is well understood.
BWA-Meth~\citep{Pedersen:2014aa} is a popular wrapper around BWA-MEM
and has been shown to have good performance in recent benchmarks~\citep{Foox:2021aa,Gong_2022,Kerns_2025,Lin_2025}.
It concatenates C-to-T and G-to-A converted reference genomes and builds one FM-index.
For paired-end reads produced with the directional protocol,
BWA-Meth converts {\tt C} bases on read 1 to {\tt T} and converts {\tt G} bases on read 2 to {\tt A}.
It aligns converted reads to the FM-index with stock BWA-MEM
and postprocesses the BWA-MEM output to generate the final alignment.
BWA-Meth has several limitations as a wrapper.
First, with the directional protocol, read 1 should be mapped either to the forward strand of the C-to-T reference or the reverse strand of the G-to-A reference.
Unaware of this constraint, BWA-MEM will also attempt to map read 1 to the reverse C-to-T and forward G-to-A reference genomes.
Although these hits are later filtered out by BWA-Meth,
they waste computing time and may affect mapping quality calculation.
Second, if read 1 is mapped to the C-to-T reference, read 2 is expected to get mapped to the G-to-A reference
--- the two reads in a pair are not located on the same contig.
This breaks BWA-MEM read pairing.
Third, BWA-MEM attempts to rescue a read if its mate is well aligned.
However, because the two reads are not supposed to be aligned to the same contig due to base conversion,
mate rescue would fail.

BISCUIT~\citep{Zhou:2024ac} and BSBolt~\citep{Farrell:2021aa} are also built upon BWA-MEM.
They avoid the issues mentioned above by customizing BWA-MEM source code specifically for BS-seq data.
They are theoretically more rigorous than BWA-Meth.
As a tradeoff, though, the deep code customization makes it challenging to adopt the faster BWA-MEM2 implementation.
Combining an efficient mapping algorithm with the BISCUIT/BSBolt strategy will yield an even better BS-seq mapper.

In this article, we will describe minibwa, the next iteration of BWA-MEM.
Minibwa aligns standard WGS short reads and Hi-C reads like BWA-MEM
and additionally maps accurate long reads and natively supports BS-seq data.

\section{Methods}

In a nutshell, minibwa is a hybrid of BWA-MEM~\citep{Li:2013aa} and minimap2~\citep{Li:2018ab},
combining BWA-MEM variable-length seeding with minimap2 chaining and SIMD-based alignment.
It additionally adapts the ropebwt3~\citep{Li:2024ac} algorithm for finding supermaximal exact matches (SMEMs).
Table~\ref{tab:relation} summarizes the relationship between minibwa and other algorithms.
BWA-MEM2 also batches suffix array (SA) query and implements SIMD-based alignment.
Further performance gains in minibwa primarily come from the new SMEM algorithm,
more frequent use of ungapped alignment and reduced effort in highly repetitive regions such as centromeres.

\begin{table}[b]
\caption{Relationship with BWA-MEM, ropebwt3 and minimap2\label{tab:relation}}
\begin{tabular*}{\columnwidth}{@{\extracolsep\fill}lll@{\extracolsep\fill}}
\toprule
Component & Source & Modification \\
\midrule
FM-index  & BWA-MEM  & Near identical \\
SMEM      & ropebwt3 & Reimplemented and batched \\
SA query  & ropebwt3 & Batched differently \\
Seeding   & BWA-MEM  & Different algorithm \\
Chaining  & minimap2 & Adapted for variable-length seeds \\
Alignment & minimap2 & More frequent ungapped fast path \\
Pairing   & BWA-MEM  & Ungapped fast path \\
\botrule
\end{tabular*}
\end{table}

For the completeness of the method section, we will describe the basic theory on FM-index in brief.
Additional technical details and explanation can be found in \citet{Li:2024ac}.
Note that the original BWA paper~\citep{Li:2009uq} used closed intervals but here we are using half-closed-half-open intervals.

\subsection{Definitions}

$\Sigma\triangleq\{{\tt A},{\tt C},{\tt G},{\tt T}\}$ is the DNA alphabet
with lexicographical order ${\tt A}<{\tt C}<{\tt G}<{\tt T}$.
$P\in\Sigma^*$ is a \emph{string} of length $|P|$.
$P[i]$, $0\le i<|P|$, is the $i$-the symbol and
$P[i,j)$ is the $(j-i)$-long substring starting at position $i$.
Operator $\circ$ concatenates two strings.
It may be omitted if string concatenation is apparent in the context.
For $a\in\Sigma$, $\overline{a}$ is the Watson-Crick complement of $a$.
$\overline{P}$ is the reverse complement of string $P$.

$R\in\Sigma^*$ is the reference genome with all contigs concatenated.
$T\triangleq R\circ\overline{R}\circ\$$ further concatenates $R$ and its reverse complement
with a trailing \emph{sentinel} $\$$ that is smaller than all symbols in $\Sigma$.
For convenice, let $n=|T|$ and $T[-1]=T[n-1]=\$$.
The \emph{suffix array} (\emph{SA}) of $T$ is an integer array $S$ where
$S(i)$ is the position of the $i$-th smallest suffix of $T$.
$B$ is the \emph{Burrows-Wheeler Transform} (\emph{BWT}) of $T$.
By definition, $B[i]=T[S(i)-1]$.

\begin{algorithm}[tb]
	\caption{Backward and forward extension}\label{algo:backfor}
	\begin{algorithmic}[1]
		\Function{BackwardExt}{$(k,k',s),a$}
			\ForAll{$b<\overline{a}$}\Comment{$b$ can be $\$$}
				\State $k'\gets k'+\big[\pi(\overline{b},k+s)-\pi(\overline{b},k)\big]$
			\EndFor
			\State $s\gets \pi(a,k+s)-\pi(a,k)$
			\State $k\gets \pi(a,k)$
			\State \Return{$(k,k',s)$}
		\EndFunction
		\Function{ForwardExt}{$(k,k',s),a$}
			\State $(k',k,s)\gets${\sc BackwardExt}$((k',k,s),\overline{a})$
			\State \Return{$(k,k',s)$}
		\EndFunction
	\end{algorithmic}
\end{algorithm}

The \emph{double-strand suffix array interval} (\emph{ds-interval}) of $P$
is a 3-tuple $({\rm lo}(P),{\rm lo}(\overline{P}),{\rm occ}(P))$ with
${\rm lo}(P)\triangleq|\{i:T[i,n)<P\}|$ and ${\rm occ}(P)\triangleq|\{i:T[i,i+|P|)=P\}|$.
If the ds-interval of $P$ is $(k,k',s)$,
function {\sc BackwardExt()} in Algorithm~\ref{algo:backfor} calculates the ds-interval of $aP$
and {\sc ForwardExt()} calculates the ds-interval of $Pa$,
where $\pi(a,k)\triangleq |\{i:B[i]<a\}|+|\{i<k:B[i]=a\}|$ is the LF-mapping.
Please see \citet{Li:2024ac} for the proof of this algorithm.

4-tuple $(P,i,j,(k,k',s))$ is a $(\ell,c)$-match
if $(k,k',s)$ is the ds-interval of $P[i,j)$, $j-i\ge\ell$ and $s\ge c$;
it is a $(\ell,c)$-SMEM if there does not exist another $(\ell,c)$-match $(P[i',j'), (\tilde{k},\tilde{k}',\tilde{s}))$ such that $i'\le i<j\le j'$.
Intuitively, a $(\ell,c)$-match involves a $P$'s substring that is $\ge\ell$ in length and occurs $\ge c$ times in $T$;
a $(\ell,c)$-SMEM is a $(\ell,c)$-match that is not contained in other $(\ell,c)$-matches on the query string.
Section~\ref{sec:smem} will describe the algorithm to find $(\ell,c)$-SMEMs.

\subsection{The memory layout of BWT}

Minibwa and BWA-MEM encode the BWT the same way.
They split $B$ into blocks.
In each block, the first four 64-bit integers store the counts of {\tt A}/{\tt C}/{\tt G}/{\tt T} ahead of the block;
the next four 64-bit integers encode 128bp sequences in $B$ in the 2-bit encoding.
Each block takes 64 bytes both on disk and in memory.
When calculating $\pi(a,k)$, minibwa finds the address of the block that contains $B[k]$,
reads the counts at the beginning of the block and scans the encoded sequences to get the exact count up to $k$.
The BWT of the human genome takes $\sim$3 GB in RAM and is too large to fit the CPU cache.
The calculation of most $\pi(a,k)$ incur cache misses.

\subsection{A brief introduction to prefetching}

\begin{figure}[tb]
\includegraphics[width=\columnwidth]{fig-prefetch.pdf}
\caption{Example of memory prefetch.
{\bf (a)} Standard backward search for three 5bp query strings.
$(r,i)$ denotes query $r$ at its position $i$.
Black arrows follow the computation order on a CPU.
Red arrows indicate data access in RAM.
{\bf (b)} Batched backward search with prefetch.}\label{fig:prefetch}
\end{figure}

A CPU takes hundreds of CPU cycles to access data in RAM but only several cycles to read data in the L1 cache.
If the data needed for computing is in RAM, the CPU has to wait for hundreds of CPU cycles before doing actual work.
Take the standard backward search as a case study.
With the classical algorithm, we process each query string in serial (Fig.~\ref{fig:prefetch}a).
Because the backward search at query position $(2,2)$ depends on the result on $(2,3)$,
the CPU has to wait for hundreds of CPU cycles to fetch data,
even though the computation itself only takes tens of CPU cycles approximately.

A more efficient approach is to decouple data access and computation by batching query strings (Fig.~\ref{fig:prefetch}b).
For example, after backward extension at $(2,3)$, we will know the address of data needed for computing $(2,2)$.
We can prefetch the data at this address, indicated by the red arrow from $(2,3)$ to $(2,2)$, and move onto position $(3,3)$.
When we come to $(2,2)$ following the black arrows, the data will be ready in cache and the CPU can perform backward extension of $(2,2)$ without waiting for the data.
This way we keep the CPU busy and effectively hide the latency of accessing data in RAM.

\subsection{Batched locate operation}

Suppose the ds-interval of $P$ is $(k,k',s)$.
It only gives positions in the SA coordinate.
To retrieve positions in $T$, we have to use suffix array $S$:
$\{S(i):k\le i<k+s\}$ is the set of $P$'s positions in $T$.
To reduce memory with FM-index, we only store $S(i)$ when $i$ is a multiple of parameter $r$.
The ``locate'' operation finds the positions of $P$ given a sampled suffix array $S$.

Function {\sc LocateSlow()} in Algorithm~\ref{algo:sa} shows the canonical implementation of the locate operation which processes each SA position in turn.
Line 8 triggers data acess and computation at the same time.
Function {\sc BatchLocate()} leverages the latency hiding technique.
Line 20 corresponds to black arrows in Fig.~\ref{fig:prefetch}b where the CPU does work;
line 21 corresponds to red arrows where prefetch happens.
Furthermore, when the input set $E$ consists of contiguous integers from a ds-interval,
line 20 will hit cache often and also greatly reduce cache misses.

To evaluate the performance of the two algorithms in Algorithm~\ref{algo:sa},
we filled $E$ with 20 random integers in range $[0,n)$.
On the FM-index of GRCh38, {\sc BatchLocate()} was over four times as fast as {\sc LocateSlow()} with identical output.
Latency hiding is an effective technique.

\begin{algorithm}[bt]
	\caption{Retrieve chromosomal positions}\label{algo:sa}
	\begin{algorithmic}[1]
		\LComment{$S$: suffix array with values only at a multiple of $r$}
		\LComment{$E$: set of suffix array positions}
		\Function{LocateSlow}{$S,r,E$}
			\State $F\gets\emptyset$
			\For{$k\in E$}
				\State $i\gets0$
				\While{$k\!\!\mod r\not=0$}
					\State $k\gets\pi(B[k],k)$
					\State $i\gets i+1$
				\EndWhile
				\State $F\gets F\cup\{i+S(k)\}$
			\EndFor
			\State \Return $F$\Comment{Set of chromosomal positions}
		\EndFunction
		\Function{BatchLocate}{$S,r,E$}
			\State $i\gets0; F\gets\emptyset$
			\While{$E\not=\emptyset$}
				\State $E'\gets\emptyset$
				\For{$k\in E$}
					\If{$k\!\!\mod r=0$}\Comment{$k$ is a multiple of $r$}
						\State $F\gets F\cup\{i+S(k)\}$
					\Else
						\State $l\gets\pi(B[k],k)$
						\State $E'\gets E'\cup\{l\}$
						\State Prefetch the memory block containing $B[l]$
					\EndIf
				\EndFor
				\State $i\gets i+1; E\gets E'$
			\EndWhile
			\State \Return $F$\Comment{Set of chromosomal positions}
		\EndFunction
	\end{algorithmic}
\end{algorithm}

\subsection{Finding SMEMs}\label{sec:smem}

Minibwa uses $(\ell,c)$-SMEMs to seed alignment.
Generalizing \citet{DBLP:conf/dlt/Gagie24} to the case of $c>1$,
\citet{Li:2024ac} provides an algorithm to find $(\ell,c)$-SMEMs given one query string.
This algorithm involves multiple rounds of forward and backward extensions (Algorithm~\ref{algo:backfor}).
For almost every extension, the CPU has to wait for the result of previous extension and then fetch the needed memory to cache.
The memory lateny is the bottleneck.

Algorithm~\ref{algo:smem} alleviates this issue by batching query strings.
Overall, it uses a queue $Q$ to organize operations on all queries in a pipeline.
Same as the original serial algorithm~\citep{Li:2024ac},
the batched version tests if $P[i,i+\ell)$ is a match with backward extension (line 7--11 in stage 1).
If it is not a match, we jump and try again (line 32);
otherwise, we use forward extension to find the end of the SMEM starting at $i$ (line 12--21 in stage 2).
A second round of backward extension is applied to find the start of the next SMEM (line 22--26 in stage 3).
Minibwa prefetches memory blocks containing $B[l]$ and $B[l']$ after line 18 and 30.
It additionally precomputes the ds-intervals of all 10-mers to reduce the number calls to {\sc BackStep}().

\begin{algorithm}[tb]
	\caption{Find $(\ell,c)$-SMEMs in string set $\mathcal{P}$}\label{algo:smem}
	\begin{algorithmic}[1]
		\Function{BatchFindSMEM}{$\mathcal{P},\ell,c$}
			\ForAll{$P\in\mathcal{P}$}\Comment{initialize queue $Q$}
				\If{$\ell-1<|P|$}
					\State $Q.{\rm push}(1,P,0,\ell-1,(0,0,|B|))$
				\EndIf
			\EndFor
			\While{$Q$ is not empty}
				\State $(u,P,i,j,(k,k',s))\gets Q.{\rm shift}()$\Comment{take out the head}
				\If{$u=1$}\Comment{stage 1: backward extension}
					\If{$j<i$}\Comment{move to stage 2}
						\State $Q.{\rm push}(2,P,i,i+\ell,(k,k',s))$
					\Else
						\State {\sc BackStep}$(Q,u,P,i,j,(k,k',s),\ell,c)$
					\EndIf
				\ElsIf{$u=2$}\Comment{stage 2: forward extension}
					\If{$j=|P|$}\Comment{reach the end of $P$}
						\State Output $(P,i,j,(k,k',s))$\Comment{report an SMEM}
						\State {\bf continue}
					\EndIf
					\State $(l,l',t)\gets${\sc ForwardExt}$((k,k',s),P[j])$
					\If{$t\ge c$}\Comment{1-step forward ext.}
						\State $Q.{\rm push}(2,P,i,j+1,(l,l',t))$
					\Else\Comment{move to stage 3}
						\State Output $(P,i,j,(k,k',s))$\Comment{report an SMEM}
						\State $Q.{\rm push}(3,P,i,j,(0,0,|B|))$
					\EndIf
				\ElsIf{$u=3$}\Comment{stage 3: backward ext. again}
					\If{$j\le i$}\Comment{back to stage 1}
						\State $Q.{\rm push}(1,P,j+1,j+\ell,(0,0,|B|))$
					\Else
						\State {\sc BackStep}$(Q,u,P,i,j,(k,k',s),\ell,c)$
					\EndIf
				\EndIf
			\EndWhile
		\EndFunction
		\Function{BackStep}{$Q,u,P,i,j,(k,k',s),\ell,c$}\Comment{helper func.}
			\State $(l,l',t)\gets${\sc BackwardExt}$((k,k',s),P[j])$
			\If{$t\ge c$}
				\State $Q.{\rm push}(u,P,i,j-1,(l,l',t))$
			\ElsIf{$j+\ell<|P|$}
				\State $Q.{\rm push}(1,P,j+1,j+\ell,(0,0,|B|))$
			\EndIf
		\EndFunction
	\end{algorithmic}
\end{algorithm}

BWA-MEM uses an old algorithm to find SMEMs~\citep{Li:2012fk}.
On real short reads without batching, the minibwa algorithm is 10\% faster to find $(19,1)$-SMEMs due to the 10-mer cache.
With batching, minibwa is $\sim$2.5 times as fast.
Minibwa is $\sim$4.6 times as fast as BWA-MEM for finding $(19,2)$-SMEMs.
The gap is larger because Algorithm~\ref{algo:smem} is more efficient to skip short SMEMs.

\subsection{Seeding, chaining and base alignment}\label{sec:align}

Minibwa seeds a query sequence $P$ in two rounds.
It first finds all $(19,1)$-SMEMs.
Suppose $(P,i,j,(k,k',s))$ is a resulting SMEM.
If $j-i\ge38$ and $s\le10$,
minibwa looks for $(\lfloor(j-1)/2\rfloor,s+1)$-SMEMs within $P[i,j)$.
During early development of BWA-MEM, we tried the same seeding algorithm.
However, as the old SMEM-finding algorithm is slow for $(19,2)$-SMEMs in the second round,
we ended up with a different seeding strategy in BWA-MEM.

The chaining algorithm in minibwa is similar to minimap2~\citep{Li:2018ab} except minor adjustment for variable-length seeds.
Minibwa retains up to 50 best chains for a short read
and uses the minimap2 algorithm to close gaps between seeds or extends the first or the last seeds.
Minibwa applies SIMD-based dynamic programming~\citep{Suzuki:2018aa} if an ungapped alignment contains many mismatches.

In comparison, BWA-MEM uses a heuristic chaining algorithm that is less accurate.
It only performs extension alignment from each seed but does not attempt to close gaps between seeds.
For reads from highly repetitive regions,
BWA-MEM tends to test more chains than minibwa but to fill a smaller portion of the dynamic programming matrix in each extension.
Overall, BWA-MEM still pays more effort for reads in centromeres or acrocentric short arms.

\subsection{Paired-end mapping}

Minibwa finds properly paired reads using the same logic as BWA-MEM.
If a read is mapped but its mate is not mapped nearby,
minibwa locally maps the mate in a window close to the mapped read.
On real data, this involves Smith-Waterman alignment~\citep{Smith:1981aa} between a 150bp read and a reference substring of several hundred basepairs in length.
Unlike BWA-MEM that always initiates full alignment,
minibwa prefilters the read-reference pair as follows.

(LATER)

\subsection{Parameter setting based on read lengths}

BWA-MEM is primarily developed for short reads.
While minimap2 supports both short and long reads,
users have to specify a preset specific for either short or long reads but not both.
This is not ideal for mixed short and long reads.
Minibwa adjusts several parameters based on the read length $\ell$ using the following formula:
$$
\theta(\ell)=\theta_l-(\theta_l-\theta_s)\cdot 2^{-\frac{\ell-\ell_{\rm min}}{\ell_{\rm mid}-\ell_{\rm min}}}
$$
where $\theta_l$ is the long-read value and $\theta_s$ is the short-read value.
We hardcode $\ell_{\rm min}=100$ and $\ell_{\rm mid}=2000$.
The length-adjusted parameter $\theta(\ell)$ reaches $(\theta_s+\theta_l)/2$ at $\ell=\ell_{\rm mid}$.

\subsection{Mapping directional bisulfite sequencing reads}

Like BWA-Meth~\citep{Pedersen:2014aa}, minibwa concatenates the C-to-T converted reference genome
and the G-to-A converted reference genome and constructs one FM-index.
This index consists of four copies of the reference genome:
forward C-to-T converted genome, forward G-to-A genome,
reverse G-to-A converted genome and reverse C-to-T genome.
Given a BS-seq read pair, minibwa converts all C bases on read 1 to T
and converts all G bases on read 2 to A.
It finds seeds against the converted FM-index.
A seed on read 1 is only retained if it comes from the forward C-to-T genome or the reverse G-to-A genome;
similarly, a seed on read 2 is retained if it comes from the forward G-to-A or the reverse C-to-T genome.
Minibwa performs base alignment and mate rescue entirely in the converted space.

\section{Results}

We evaluated the mapping accuracy of minibwa on simulated short and long reads and
measured the performance on real data including WGS reads,
Hi-C short reads and directional BS-seq reads.

\subsection{Mapping accuracy on simulated data}

\begin{figure*}
\includegraphics[width=\textwidth]{fig-roc.pdf}
\caption{Mapping accuracy measured on simulated reads.
Short and long reads are simulated from GRCh38, T2T-CHM13 or the HG002 diploid genomes and are mapped to the corresponding source genomes with various read aligners.
Suppose a read is originated from genomic interval $G_s$ in simulation
and is mapped to interval $G_a$ by an aligner.
It is considered mapped correctly if $|G_s\cap G_a|/|G_s\cup G_a|\ge10\%$.
Each point in the figures gives the fraction of mapped reads and the fraction of correctly mapped at a mapping quality threshold.
{\bf (a)} $2\times150$ bp paired-end reads simulated from GRCh38.
{\bf (b)} $2\times150$ bp paired-end reads simulated from T2T-GRCh38.
{\bf (c)} Same reads in {\bf (b)} but excluding reads from centromeres or acrocentric short arms.
{\bf (d)} Same reads in {\bf (a)} but mapped in the methylation mode.
{\bf (e)} HiFi and Nanopore reads simulated from T2T-CHM13.
{\bf (f)} HiFi reads simulated from the diploid HG002 genome consisting of two haplotypes.
}\label{fig:roc}
\end{figure*}

We simulated short reads using mason~\citep{fu_mi_publications962} with
``{\tt -{}-illumina-read-length 150 -{}-illumina-prob-mismatch-scale\\2.5}''.
We simulated long reads with Badread~\citep{Wick2019},
with the default option for Nanopore reads
and option ``{\tt -{}-error\_model pacbio2021 -{}-qscore\_model pacbio2021 -{}-length 15000,3000\\-{}-identity 27,3}'' for HiFi reads.
We encoded the true read coordiates in read names with paftools in the minimap2 package.
and plotted the results (Fig.~\ref{fig:roc}).
A mapper of higher accuracy tends to yield a curve closer to the upper left corner.

BWA-MEM is slightly more accurate than minibwa for WGS short reads (Fig.~\ref{fig:roc}a and Fig.~\ref{fig:roc}b).
The difference is primarily driven by centromeric or acrocentric regions (Fig.~\ref{fig:roc}b versus Fig.~\ref{fig:roc}c)
as minibwa pays less effort in such regions (Section~\ref{sec:align}).
Minibwa competes with minimap2~\citep{Li:2018ab} and Winnowmap2~\citep{Jain:2020aa,Jain:2022ua} for mapping accurate long reads (Fig.~\ref{fig:roc}e and Fig.~\ref{fig:roc}f).

All aligners shown in Fig.~\ref{fig:roc}d use the 3-base strategy and should be able to align normal WGS reads.
Both BWA-Meth~\citep{Pedersen:2014aa} and BISCUIT~\citep{Zhou:2024ac} are built on top of BWA-MEM.
BWA-Meth is not as accurate as minibwa due to its limitations explained in Introduction.
In theory, BISCUIT uses a better strategy than minibwa because
it properly penalizes a mismatch between a T base on the reference and a C on the read --- this would not happen unless there were a real SNP.
It is not clear to us why BISCUIT is less accurate than minibwa.
Bismark is known to map fewer reads than BWA-Meth~\citep{Foox:2021aa,Gong_2022,Kerns_2025,Lin_2025}.

We note that on real data, most false read alignment at high mapping quality is caused by large structural variants between individuals.
For example, if the reference genome has one copy of a gene but the sequenced sample has an extra copy,
all reads from the extra copy will be mismapped at high mapping quality.
The theoretical accuracy as is shown in Fig.~\ref{fig:roc} may not be strongly correlated with downstream analyses.

\subsection{Mapping performance on real data}

\begin{figure*}
\includegraphics[width=\textwidth]{fig-perf}
\caption{Speed and peak memory on real data.}\label{fig:perf}
\end{figure*}

\subsection{Effect on short-read variant calling}

\section{Discussions}

\section*{Acknowledgments}

blabla

\section*{Author contributions}

H.L. conceived the project, implemented the core algorithm, analyzed the data and drafted the manuscript.
N.H. added new features, contributed to development and reviewed the draft.

\section*{Conflict of interest}

N.H. is an employee of Fulcrum Genomics LLC.

\section*{Funding}

This work is supported by National Institute of Health grant R01HG010040, R01HG014175 and U24CA294203 (to H.L.).

\section*{Data availability}

The minibwa source code is available at \url{https://github.com/lh3/minibwa}.

\bibliographystyle{apalike}
{\sffamily\small
\bibliography{minibwa}}

\end{document}
